%% kapa: exam=1 round=32 year=2021 period=1 form=A subject=economics subject_ko=경제학원론 seq=33 number=73
%% source: Q-Net 감정평가사 기출문제 (cst00309) artlSeq=5208397 file=extracted/1008/5208397/A형/2021년도 제32회 감정평가사 1차 1교시 A형.pdf
%% conversion: kordoc 4.24.51 -> tools/gen1.py (auto); review=proofread-sonnet; flags=none
\begin{question}{73}
현재 인플레이션율 8\%에서 4\%로 낮출 경우, 보기를 참고하여 계산된 희생률은? (단, $\Pi_t$, $\Pi_{t-1}$, $U_t$는 각각 t기의 인플레이션율, (t-1)기의 인플레이션율, t기의 실업률이다.)
\begin{bogi}
○ $\Pi_t-\Pi_{t-1}=-0.8(U_t-0.05)$ \par
○ 현재 실업률: 5\% \par
○ 실업률 1\%p 증가할 때 GDP 2\% 감소로 가정 \par
○ 희생률: 인플레이션율을 1\%p 낮출 경우 감소되는 GDP 변화율(\%) \par
\end{bogi}
\begin{choices}
  \item 1.5
  \item 2
  \item 2.5
  \item 3
  \item 3.5
\end{choices}
\end{question}
\answer{3}
